\documentclass[11pt]{article}
\usepackage[margin=1in]{geometry}
\usepackage{amsmath,amssymb,amsthm,booktabs,array,hyperref}
\hypersetup{colorlinks=true,urlcolor=blue,citecolor=blue,linkcolor=blue}
\newtheorem{theorem}{Theorem}
\newtheorem{proposition}[theorem]{Proposition}
\newcommand{\D}{D}
\newcommand{\code}[1]{\texttt{\small #1}}
\setlength{\parskip}{5pt}
\setlength{\parindent}{0pt}
\title{Heterogeneous parts in fractional defect:\\an eight-vertex counterexample}
\author{Research draft for MathDB 335230}
\date{1 November 2026 \quad Draft only; literature clearance on hold}
\begin{document}
\maketitle
\begin{abstract}
An explicit three-color usage assignment on the complete four-partite graph
with two vertices in each part has maximum fractional defect $5/3$.
Consequently $D(K_2^{(4)},3)\leq 5/3<2$, refuting the proposed universal
formula $D(K_a^{(m)},k)=(\lceil m/k\rceil-1)a$ in the inspected sources.
The construction uses different usages inside two independent parts; every
coloring constant on each part has maximum defect at least $2$.
An unchanged Lean certificate proves the upper bound and both minimum and
infimum formulations of the universal refutation. The optimum is unknown
here. A directly relevant 2016 thesis body remains inaccessible, so this
draft makes no literature-priority claim and carries no acceptance credit.
\end{abstract}

\section{Definitions and the exact target}
For a finite nonempty simple graph $G$ and a positive integer $k$, a usage assignment is
$f:V(G)\to\Delta_{k-1}$, where
$\Delta_{k-1}=\{(x_1,\ldots,x_k)\in\mathbb R^k:x_j\geq0,\ \sum_jx_j=1\}$.
Define overlap, vertex defect, and maximum defect by
\begin{align*}
 o_f(v,w)&=\sum_{j=1}^k\min\{f_j(v),f_j(w)\},\\
 d_f(v)&=\sum_{w\in N(v)}o_f(v,w),&
 \delta(f)&=\max_{v\in V(G)}d_f(v).
\end{align*}
The invariant $D(G,k)=\min_f\delta(f)$ exists: the feasible set is a finite
product of compact simplexes and the objective is continuous. This is the
source definition \cite[Section 2, equation (1)]{v2}, not the usual
fractional chromatic number. Usages are independent at every vertex.

For positive integers $m,a$, let $K_a^{(m)}$ have $m$ independent parts of size
$a$, with every edge between distinct parts. MathDB 335230 asks whether,
for \emph{all positive integers} $m,a,k$,
\begin{equation}\label{conj}
 D(K_a^{(m)},k)=(\lceil m/k\rceil-1)a.
\end{equation}
This is Conjecture 2 in both arXiv versions \cite{v1,v2} and Conjecture 1
in the undated author preprint \cite{author}. Section 3's default two-color
convention does not restrict this explicitly $k$-color conjecture.
The frozen literal record and its corpus-line provenance are included in the dossier;
its historical open label and empty solution list are not novelty evidence.

\newpage
\section{The witness and universal refutation}
\begin{theorem}\label{main}
For the source definition above, $D(K_2^{(4)},3)\leq5/3<2$.
In particular, the universal equality \eqref{conj} is false.
\end{theorem}
\begin{proof}
Name the parts $A,B,C,D$, each with vertices indexed by $0,1$. Assign the
following usages, with coordinates in the order colors $1,2,3$:
\begin{center}
\begin{tabular}{c c c}
\toprule Vertex & Usage vector & Vertex defect \\ \midrule
$A_0,A_1$ & $(0,0,1)$ & $2/3$ each\\
$B_0,B_1$ & $(0,2/3,1/3)$ & $5/3$ each\\
$C_0$ & $(0,1,0)$ & $5/3$\\
$C_1$ & $(1,0,0)$ & $5/3$\\
$D_0$ & $(2/3,1/3,0)$ & $5/3$\\
$D_1$ & $(1,0,0)$ & $1$\\ \bottomrule
\end{tabular}
\end{center}
Every vector is nonnegative with coordinate sum one. Here is the complete
accounting of the nonzero overlaps on \emph{edges}; symmetry supplies the
reverse ordered pairs:
\[
\begin{array}{c|c|c}
\text{edge type}&\text{indices}&\text{overlap}\\ \hline
A_iB_j&i,j\in\{0,1\}&1/3\\
B_iC_0&i\in\{0,1\}&2/3\\
B_iD_0&i\in\{0,1\}&1/3\\
C_0D_0&&1/3\\
C_1D_0&&2/3\\
C_1D_1&&1
\end{array}
\]
All other cross-part overlaps are zero. Within-part pairs are not edges,
regardless of their usages. For example
\[
 o_f(B_0,D_0)=\min(0,2/3)+\min(2/3,1/3)+\min(1/3,0)=1/3.
\]
Thus $d_f(A_i)=2(1/3)=2/3$ and
$d_f(B_i)=2(1/3)+2/3+1/3=5/3$.
Also $d_f(C_0)=2(2/3)+1/3=5/3$,
$d_f(C_1)=2/3+1=5/3$,
$d_f(D_0)=2(1/3)+1/3+2/3=5/3$, and $d_f(D_1)=1$.
The exhibited coloring therefore has $\delta(f)=5/3$.
Taking the minimum gives $D\leq5/3<2$.
But the proposed value at the positive triple $(m,a,k)=(4,2,3)$ is
$(\lceil4/3\rceil-1)2=2$. A universal equality would hold at this triple,
contradicting the strict inequality.
\end{proof}
The equality $\delta(f)=5/3$ concerns this witness alone. Neither
$D(K_2^{(4)},3)=5/3$ nor minimal counterexample size is asserted.

\newpage
\section{Why heterogeneity matters; comparison with known results}
\begin{proposition}
Among usage assignments on $K_2^{(4)}$ constant within each part, the least
maximum defect for three colors is $2$. Among pure three-color assignments
on this graph, the least maximum defect is also $2$.
\end{proposition}
\begin{proof}
A constant vector on each part gives a three-color usage assignment $g$ on
$K_4$. For every vertex in a part, each of the other three vectors is
encountered twice, so $\delta(f)=2\delta(g)$.
The source complete-graph theorem gives $D(K_4,3)=1$
\cite[Proposition 1]{v2}; its replicated optimum attains $2$.

For the pure case, one color occurs at least three times among eight vertices.
Choose three occurrences. Since a part has size two, if two lie in one part,
the third has two same-color neighbors; if they lie in three parts, each has
two such neighbors. Hence the maximum is at least $2$. Giving $A$ and $B$
one color and $C,D$ distinct other colors attains $2$.
\end{proof}
The witness strictly leaves both classes. The parts $C,D$ are heterogeneous,
even though their vertices have identical neighborhoods. Thus independent
twins cannot in general be constrained to one replicated vector without
changing the optimum. With $G=K_4,a=2,k=3$, it also disproves the broader
replication equality suggested immediately after the source conjecture;
Proposition 4(b)'s upper inequality itself remains valid.

The following comparisons use the actual statements and relevant proof
bodies of \cite{v1,v2,author}, rather than just their abstracts:
\begin{itemize}
\item The complete-graph result is
$D(K_n,k)=\lceil n/k\rceil-1$. Combined with replication it gives the upper
bound $2$ here, and the preceding proposition shows it cannot give a strict
improvement within that construction.
\item Proposition 4(a) and Corollary 1 of \cite{v2} give
$TD(K_2^{(4)},3)\geq4TD(K_4,3)=8$, where $TD$ is minimum total defect.
Averaging gives only $D\geq1$. The monochromatization lemma for total defect
need not preserve the maximum defect of individual vertices.
\item Proposition 5 of \cite{v2} proves \eqref{conj} when $k\mid m$;
$3\nmid4$. A six-part example with $a=2,k=3$ has bound $2$;
restricting it to four parts supplies no strict bound from that argument.
\item Proposition 6 proves \eqref{conj} for $k=2$. Its proof orders a scalar
red usage $r$ and uses overlap $1-|r(u)-r(v)|$ with blue usage $1-r$.
This argument is not a common scalar ordering of three-coordinate usages.
Reserving a third pure color for one whole part and applying that theorem
to the remaining three parts again gives $2$.
\end{itemize}
These are specific limitations of the inspected constructions, not a theorem
excluding every possible deduction from the literature. The strict separation
is a substantive candidate mechanism; the short witness and formalization
alone do not establish novelty.

\newpage
\section{Lean correspondence and logical endpoints}
The dossier preserves \code{FractionalDefect.lean} byte for byte. All names
below are in namespace \code{FractionalDefect}. Natural indices $m,a,k$
encode the source parameters $M=m+1,A=a+1,K=k+1$, so the witness's indices
are $(3,1,2)$, not a zero-parameter exception.
\begin{center}\small
\begin{tabular}{p{0.27\linewidth}p{0.64\linewidth}}
\toprule Source notion & Lean declaration and interpretation\\\midrule
Vertices and edges & \code{Vertex} is $\mathrm{Fin}(m+1)\times\mathrm{Fin}(a+1)$;
\code{Adj} is inequality of first coordinates.\\
Independent usages & \code{Coloring} assigns real coordinates;
\code{Valid} asserts nonnegativity and unit sum at each vertex.\\
Objective & \code{defect} sums overlap on edges, zero on nonedges;
\code{graphDefect} is the nonempty finite supremum.\\
Claimed minimum & \code{MinimumDefect} is \code{IsLeast} of exactly the
valid attainable objective values.\\
Numerical invariant & \code{D} is the real infimum of the same set;
\code{D\_eq\_of\_minimum} identifies an attained minimum with it.\\
Conjecture value & \code{predicted} uses natural quotient
$(M+K-1)/K$, cast to $\mathbb R$, then subtracts $1$ and multiplies by $A$.\\
Universal statements & \raggedright\arraybackslash\code{SourceConjecture} uses \code{MinimumDefect};
\code{NumericalSourceConjecture} uses equality for \code{D}.\\\bottomrule
\end{tabular}\end{center}
For completeness, the ceiling encoding is justified independently: write
$M=qK+r$, $0\leq r<K$. If $r=0$, both $\lceil M/K\rceil$ and
$\lfloor(M+K-1)/K\rfloor$ equal $q$; otherwise both equal $q+1$.
This is a semantic bridge in this paper, not a claim that the file contains
a generic real-ceiling theorem. The relevant value $2$ is kernel checked.

\code{witness\_valid} proves all simplex constraints by finite cases and exact
rational normalization. \code{witness\_defects} proves the upper bound at
each vertex (its name does not assert the exact row values).
\code{graphDefect\_witness\_le} lifts this bound to the maximum.
The exact table above also has a separate rational-arithmetic receipt.

\code{defect\_nonneg} and \code{graphDefect\_nonneg} bound the attainable
set below by zero. The witness supplies a member, so
\code{D\_witness\_le} applies the bounded-infimum inequality and proves
$D\ 3\ 1\ 2\leq5/3$.
\code{D\_strict\_counterexample} proves strict inequality against
\code{predicted}. \code{literal\_minimum\_refuted} proves
$\neg\,\code{MinimumDefect}\ 3\ 1\ 2\ 2$: any least value $2$ would be
at most the witness objective, which is at most $5/3$.
\code{source\_conjecture\_false} and
\code{numerical\_source\_conjecture\_false} specialize the respective
universal assertions to this triple and derive contradiction.

General minimum attainment is justified above in ordinary mathematics;
a generic Lean compactness theorem is not part of this packet and is not
needed for either refutation. Eleven mathematical theorem declarations and
all 32 local declarations, including generated helpers, are audited in a
separate unchanged review file. Only \code{propext}, \code{Classical.choice},
and \code{Quot.sound} occur. No new formal method is claimed.

\newpage
\section{Reproducibility, sources, and unresolved literature status}
\small\raggedright\setlength{\parskip}{3pt}
The original seven-file packet is preserved as \code{original/proof-packet.zip};
the outer dossier is nonrecursive. In a fresh dossier extraction run
\code{python3 verify-manifest.py}, then \code{python3 replay-dossier.py}.
The latter safely extracts the original packet to a new work directory, runs
its unchanged \code{replay.py}, and separately runs the unchanged
\code{AuditAll.lean}. It needs Python 3, Linux affinity/resource APIs, and
the pinned external Lean and dependency caches. Set \code{LEAN\_BINARY} and
\code{MATHLIB\_PACKAGES} if their recorded absolute defaults differ.
Remove only the generated \code{work/strict-extracted} directory before a
repeat run. The archive contains no compiled oleans, caches or rendered pages.

The toolchain pins are:\\
Lean 4.34.1: \code{5045d0056413266e57c625dcd7c365b10e377c52}\\
mathlib: \code{d13f23b723b8a846827a245b89c10fc7d3f11612}.
The lake manifest records nine dependency revisions. Runs use strict warnings,
\code{-j 1 -M 2048}, at most two CPUs, an 8 GiB address cap, a 90-second soft
CPU cap (95 hard), and a 105-second wall timeout. Logs and resource receipts
are bundled. This is an environment-pinned replay: prebuilt imported oleans
and their exported axiom summaries remain external trusted inputs. Revision
metadata alignment does not prove source cleanliness or source-to-olean
provenance. No full dependency rebuild or fully vendored offline environment
is claimed.

Source bodies inspected are arXiv v1 (5 February 2017), v2 (7 November 2019),
and the undated author preprint. The first two label the target Conjecture 2,
the third Conjecture 1; the definitions and target agree. The author copy has
different numbering and two-color proof organization. It is not authenticated
as a version of record. Dated arXiv metadata identifies no journal reference;
no separate journal body is known here, and none is invented as a requirement.
The author slides restrict their parameter study to two colors. The related
series-parallel paper \cite{sp} uses a different defective set-coloring model
and excludes $K_4$ minors; the present graph contains $K_4$. The review's
bounded comparisons found no direct strict-bound construction there.

\textbf{Material HOLD.} The target paper acknowledges that the work is part
of Xu's 2016 thesis \cite{xu}. Its body has not been inspected. Official body
access returned 403 around 07:14 UTC on 1 November 2026; a later landing-page
request returned 403 at 07:21:41 UTC. Existing bounded searches found no
independent lawful full copy. No denied route was retried or bypassed for this
draft, and no new independent lead arose. The original packet's receipts
predate the last body failure; the outer source log and review receipts
supplement them without changing their bytes. A later retained conjecture
does not rule out a subsuming construction in the thesis.

Independent source/ordinary and formal-fidelity reviews PASS; significance
is qualified PASS, with literature HOLD retained. Dated registry screens and
inherited MathDB/Palomar denials give no exhaustive priority clearance.
Candidate credit remains zero. Root paper review, final replay, literature
adjudication and acceptance remain separate.

\begin{thebibliography}{9}\footnotesize
\bibitem{v2} W. Goddard and H. Xu, \emph{Colorings with Fractional Defect},
arXiv:1702.01442v2, 7 Nov. 2019. Sections 2 and 3.3, Proposition 1,
Propositions 4--6 and Conjecture 2.
\url{https://arxiv.org/pdf/1702.01442v2}.
\bibitem{v1} Same authors and title, arXiv:1702.01442v1, 5 Feb. 2017.
\url{https://arxiv.org/pdf/1702.01442v1}.
\bibitem{author} Same authors and title, undated author-hosted manuscript,
Section 3.3, Conjecture 1.
\url{https://people.computing.clemson.edu/~goddard/papers/fractionalDefectPreprint.pdf}.
\bibitem{sp} W. Goddard and H. Xu, \emph{Fractional, Circular, and Defective
Coloring of Series-Parallel Graphs}, author manuscript; the inspected CV lists
J. Graph Theory \textbf{81} (2016), 146--153.
\url{https://people.computing.clemson.edu/~goddard/papers/seriesParallel.pdf}.
\bibitem{xu} H. Xu, \emph{Generalized colorings of graphs}, Clemson dissertation (2016), repository item 1650.
Body unavailable in this investigation; cited as an unresolved predecessor.
\url{https://open.clemson.edu/all_dissertations/1650/}.
\end{thebibliography}
\end{document}
